%%%%%%%%%%%%%%%%%%%%%%%%%%%%%%%%%%%%%%%%%
% fphw Assignment
% LaTeX Template
% Version 1.0 (27/04/2019)
%
% This template originates from:
% https://www.LaTeXTemplates.com
%
% Authors:
% Class by Felipe Portales-Oliva (f.portales.oliva@gmail.com) with template 
% content and modifications by Vel (vel@LaTeXTemplates.com)
%
% Template (this file) License:
% CC BY-NC-SA 3.0 (http://creativecommons.org/licenses/by-nc-sa/3.0/)
%
%%%%%%%%%%%%%%%%%%%%%%%%%%%%%%%%%%%%%%%%%

%----------------------------------------------------------------------------------------
%	PACKAGES AND OTHER DOCUMENT CONFIGURATIONS
%----------------------------------------------------------------------------------------

\documentclass[
	12pt,
	%fleqn, % Default font size, values between 10pt-12pt are allowed
	%letterpaper, % Uncomment for US letter paper size
	%spanish, % Uncomment for Spanish
]{fphw}

% Template-specific packages
\usepackage[utf8]{inputenc} % Required for inputting international characters
\usepackage[T1]{fontenc} % Output font encoding for international characters
\usepackage{mathpazo} % Use the Palatino font

\usepackage{graphicx} % Required for including images

\usepackage{booktabs} % Required for better horizontal rules in tables

\usepackage{listings} % Required for insertion of code

\usepackage{enumerate} % To modify the enumerate environment

\usepackage{enumitem}

\usepackage{amsmath}

\usepackage{mathtools,amsthm}

\usepackage{float}

\usepackage{listings}

\usepackage{hyperref}

\lstdefinelanguage{R}{
  keywords={if, else, repeat, while, function, for, in, next, break},
  otherkeywords={TRUE, FALSE, NULL, NA, Inf, NaN},
  sensitive=true,
  morecomment=[l]\#,
  morestring=[b]",
}

\newcommand{\E}{\mathbb{E}}
\newcommand{\Var}{\operatorname{Var}}
\newcommand{\R}{\mathbb{R}}
\newcommand{\norm}[1]{\left\lVert #1 \right\rVert}


%----------------------------------------------------------------------------------------
%	ASSIGNMENT INFORMATION
%----------------------------------------------------------------------------------------

\title{Homework \#1} % Assignment title

\author{Shizhe Zhang} % Student name

\date{\today} % Due date

\institute{University of California, Berkeley \\ Department of Statistics} % Institute or school name

\class{EECS 227AT} % Course or class name

\professor{Gireeja Ranade} % Professor or teacher in charge of the assignment

%----------------------------------------------------------------------------------------

\begin{document}

\maketitle % Output the assignment title, created automatically using the information in the custom commands above

%----------------------------------------------------------------------------------------
%	ASSIGNMENT CONTENT
%----------------------------------------------------------------------------------------

\section*{Problem 1. Optimizing Over Multiple Variables}


In this exercise, we consider several problems in which we optimize over two variables, $\vec{x} \in \mathbb{R}^n$ and $\vec{y} \in \mathbb{R}^m$, and a general (possibly nonconvex) objective function $F_0(\vec{x}, \vec{y})$. Suppose also that $\vec{x}$ and $\vec{y}$ are constrained to different feasible sets $X$ and $Y$, respectively, which may or may not be convex.

\begin{enumerate}
\item[(a)] Show that
\[
\min_{\vec{x}\in X}\min_{\vec{y}\in Y} F_0(\vec{x},\vec{y})
=
\min_{\vec{y}\in Y}\min_{\vec{x}\in X} F_0(\vec{x},\vec{y}),
\]
i.e., if we minimize over both $\vec{x}$ and $\vec{y}$, then we can exchange the minimization order without altering the optimal value.

\item[(b)] Show that $p^\star \ge d^\star$, where
\[
p^\star := \min_{\vec{x}\in X}\max_{\vec{y}\in Y} F_0(\vec{x},\vec{y})
\]

\[
d^\star := \max_{\vec{y}\in Y}\min_{\vec{x}\in X} F_0(\vec{x},\vec{y})
\]

This statement is referred to as the min-max theorem.
\end{enumerate}

\subsection*{Answer}

\textbf{(a)}

Let
\[
a = \min_{\vec{x}\in X}\min_{\vec{y}\in Y} F_0(\vec{x},\vec{y}),
\quad
b = \min_{\vec{y}\in Y}\min_{\vec{x}\in X} F_0(\vec{x},\vec{y}).
\]

Both expressions represent the minimum value of the same function over the same joint feasible set:
\[
(\vec{x},\vec{y}) \in X \times Y .
\]

Indeed,
\[
\min_{\vec{x}\in X}\min_{\vec{y}\in Y} F_0(\vec{x},\vec{y})
=
\min_{(\vec{x},\vec{y})\in X\times Y} F_0(\vec{x},\vec{y}).
\]

Similarly,
\[
\min_{\vec{y}\in Y}\min_{\vec{x}\in X} F_0(\vec{x},\vec{y})
=
\min_{(\vec{x},\vec{y})\in X\times Y} F_0(\vec{x},\vec{y}).
\]

Thus the order of minimization does not affect the optimal value.

\bigskip

\textbf{(b)}

Define
\[
p^\star = \min_{\vec{x}\in X}\max_{\vec{y}\in Y} F_0(\vec{x},\vec{y}),
\quad
d^\star = \max_{\vec{y}\in Y}\min_{\vec{x}\in X} F_0(\vec{x},\vec{y}).
\]

For any $\vec{x}\in X$ and $\vec{y}\in Y$, we have
\[
\min_{\vec{x}'\in X} F_0(\vec{x}',\vec{y}) \le F_0(\vec{x},\vec{y})
\le \max_{\vec{y}'\in Y} F_0(\vec{x},\vec{y}').
\]

Taking the maximum over $\vec{y}$ on the left inequality gives
\[
\max_{\vec{y}\in Y} \min_{\vec{x}\in X} F_0(\vec{x},\vec{y})
\le
\max_{\vec{y}\in Y} F_0(\vec{x},\vec{y})
\]

for every $\vec{x}$.

Now take the minimum over $\vec{x}$ on the right-hand side:

\[
\max_{\vec{y}\in Y} \min_{\vec{x}\in X} F_0(\vec{x},\vec{y})
\le
\min_{\vec{x}\in X} \max_{\vec{y}\in Y} F_0(\vec{x},\vec{y}).
\]

Therefore
\[
d^\star \le p^\star.
\]

%----------------------------------------------------------------------------------------

\newpage
\section*{Problem 2. Duality}


Consider the function
\[
f(\vec{x}) = \vec{x}^T A \vec{x} - 2\vec{b}^T \vec{x}.
\]

First, we consider the unconstrained optimization problem
\[
p^\star = \min_{\vec{x}\in\mathbb{R}^n} f(\vec{x})
= \min_{\vec{x}\in\mathbb{R}^n} \vec{x}^T A \vec{x} - 2\vec{b}^T \vec{x}
\]

for a real $n\times n$ symmetric matrix $A\in S^n$ and $\vec{b}\in\mathbb{R}^n$.  
If the problem is unbounded below, then we say $p^\star = -\infty$.  
Let $\vec{x}^\star$ denote the minimizing argument.

\begin{enumerate}
\item[(a)] Suppose $A\succeq 0$ and $\vec{b}\in \mathcal{R}(A)$. Let $\mathrm{rank}(A)=n$. Find $p^\star$.

Hint: What does $A\succeq0$ tell you about the function $f$? How can you leverage the rank of $A$ to compute $p^\star$?

\item[(b)] Suppose $A\succeq0$ and $\vec{b}\in R(A)$ as before. Let $A$ be rank-deficient, i.e., $\mathrm{rank}(A)=r<n$. Let $A$ have compact/thin and full SVD

\[
A = U_r\Lambda_rU_r^T =
\begin{bmatrix}
U_r & U_1
\end{bmatrix}
\begin{bmatrix}
\Lambda_r & 0 \\
0 & 0
\end{bmatrix}
\begin{bmatrix}
U_r^T \\
U_1^T
\end{bmatrix}
\]

Show that the minimizer $\vec{x}^\star$ is not unique by finding the family of solutions in terms of $U_r, U_1, \Lambda_r, \vec{b}$.

Hint: Any $\vec{x}\in\mathbb{R}^n$ can be written as
\[
\vec{x}=U_r\vec{\alpha}+U_1\vec{\beta}.
\]

\item[(c)] If $A\not\succeq0$, show that $p^\star=-\infty$ by finding $\vec{v}$ such that $f(\alpha\vec{v})\to -\infty$ as $\alpha\to\infty$.

Hint: $A\not\succeq0$ means there exists $\vec{v}$ such that $\vec{v}^TA\vec{v}<0$.

\item[(d)] Suppose $A\succeq0$ and $\vec{b}\notin R(A)$. Find $p^\star$.

Hint: From FTLA,
\[
\mathbb{R}^n = R(A^T) \oplus N(A).
\]

\item[(e)] Consider the constrained problem

\[
p^\star = \min_{\vec{x}\in\mathbb{R}^n} \vec{x}^T A \vec{x} - 2\vec{b}^T \vec{x}
\]

subject to

\[
\vec{x}^T\vec{x} \ge 1.
\]

Write the Lagrangian $L(\vec{x},\lambda)$.

\item[(f)] For any matrix $C\in\mathbb{R}^{n\times n}$ with compact SVD

\[
C = U_r\Lambda_rV_r^T
\]

define the pseudoinverse

\[
C^\dagger = V_r\Lambda_r^{-1}U_r^T.
\]

Show that the dual problem to (7) can be written as

\[
d^\star = \max_{\lambda\ge0}
\]

subject to

\[
A-\lambda I\succeq0,\quad \vec{b}\in R(A-\lambda I)
\]

\[
-b^T(A-\lambda I)^\dagger b + \lambda.
\]

Hint: Compute $g(\lambda)$ and analyze its behavior under the constraints.
\end{enumerate}

\subsection*{Answer}

\textbf{(a)}

Suppose $A \succeq 0$ and $\mathrm{rank}(A)=n$. Then $A$ is positive definite and invertible.

The objective is
\[
f(\vec{x}) = \vec{x}^T A \vec{x} - 2\vec{b}^T \vec{x}.
\]

Set the gradient to zero:

\[
\nabla f(\vec{x}) = 2A\vec{x} - 2\vec{b} = 0
\]

\[
A\vec{x} = \vec{b}.
\]

Since $A$ is invertible,

\[
\vec{x}^\star = A^{-1}\vec{b}.
\]

Substitute into the objective:

\[
p^\star
= (A^{-1}\vec{b})^T A (A^{-1}\vec{b}) - 2\vec{b}^T(A^{-1}\vec{b})
\]

\[
= \vec{b}^T A^{-1}\vec{b} - 2\vec{b}^T A^{-1}\vec{b}
\]

\[
p^\star = -\vec{b}^T A^{-1}\vec{b}.
\]

\textbf{(b)}

Let $A$ have compact SVD

\[
A = U_r \Lambda_r U_r^T,
\]

with $\mathrm{rank}(A)=r<n$.

Since $\vec{b}\in R(A)$, we can write

\[
\vec{b} = U_r \vec{c}
\]

for some vector $\vec{c}$.

Write any $\vec{x}\in\mathbb{R}^n$ as

\[
\vec{x} = U_r \vec{\alpha} + U_1 \vec{\beta}.
\]

Then

\[
f(\vec{x})
= (U_r\vec{\alpha}+U_1\vec{\beta})^T
A
(U_r\vec{\alpha}+U_1\vec{\beta})
-2\vec{b}^T(U_r\vec{\alpha}+U_1\vec{\beta}).
\]

Since $AU_1=0$ and $U_1^T U_r=0$,

\[
f(\vec{x})
=
\vec{\alpha}^T \Lambda_r \vec{\alpha}
-2\vec{c}^T\vec{\alpha}.
\]

Thus the objective depends only on $\vec{\alpha}$.

Minimizing gives

\[
\vec{\alpha}^\star = \Lambda_r^{-1}\vec{c}.
\]

Hence

\[
\vec{x}^\star
=
U_r \Lambda_r^{-1} \vec{c}
+
U_1 \vec{\beta}
\]

Since $\vec{c}=U_r^T\vec{b}$,

\[
\vec{x}^\star = U_r \Lambda_r^{-1} U_r^T \vec{b} + U_1 \vec{\beta},\quad \forall \vec{\beta}
\]

which shows the minimizer is not unique.

\bigskip

\textbf{(c)}

If $A \not\succeq 0$, then there exists $\vec{v}$ such that

\[
\vec{v}^T A \vec{v} < 0.
\]

Consider $\vec{x}=\alpha \vec{v}$.

Then

\[
f(\alpha\vec{v})
=
\alpha^2 \vec{v}^T A \vec{v}
-2\alpha \vec{b}^T \vec{v}.
\]

The dominant term as $\alpha\to\infty$ is

\[
\alpha^2 \vec{v}^T A \vec{v}.
\]

Since $\vec{v}^T A \vec{v}<0$,

\[
f(\alpha\vec{v}) \to -\infty.
\]

\textbf{(d)}

Suppose $A\succeq0$ but $\vec{b}\notin R(A)$.

Decompose

\[
\vec{b} = \vec{v}_1 + \vec{v}_2
\]

where

\[
\vec{v}_1 \in R(A), \quad \vec{v}_2 \in N(A).
\]

Choose $\vec{x}=t\vec{v}_2$.


\[
f(t\vec{v}_2)
=
t^2 \vec{v}_2^T A \vec{v}_2
-
2t \vec{b}^T \vec{v}_2.
\]

Since $A\vec{v}_2=0$,

\[
f(t\vec{v}_2) = -2t \vec{b}^T \vec{v}_2.
\]

But

\[
\vec{b}^T\vec{v}_2 = \vec{v}_2^T\vec{v}_2 >0.
\]

Thus

\[
f(t\vec{v}_2) \to -\infty
\]

as $t\to\infty$.
Therefore
$ p^\star = -\infty. $

\bigskip

\textbf{(e)}

The constraint is

\[
\vec{x}^T\vec{x} \ge 1.
\]

Rewrite as

\[
1-\vec{x}^T\vec{x} \le 0.
\]

The Lagrangian is

\[
L(\vec{x},\lambda)
=
\vec{x}^T A \vec{x}
-2\vec{b}^T\vec{x}
+\lambda(1-\vec{x}^T\vec{x})
\]

\[
=
\vec{x}^T(A-\lambda I)\vec{x}
-2\vec{b}^T\vec{x}
+\lambda.
\]

\bigskip

\textbf{(f)}

The dual function is

\[
g(\lambda)=\min_{\vec{x}} L(\vec{x},\lambda).
\]

If $A-\lambda I \not\succeq 0$, the quadratic is unbounded below, so

\[
g(\lambda)=-\infty.
\]

If $A-\lambda I\succeq0$, set the gradient to zero:

\[
2(A-\lambda I)\vec{x}-2\vec{b}=0.
\]

Thus

\[
(A-\lambda I)\vec{x}=\vec{b}.
\]

A solution exists only if

\[
\vec{b}\in R(A-\lambda I).
\]

The minimizing solution is

\[
\vec{x}=(A-\lambda I)^\dagger \vec{b}.
\]

Substitute into $L$:

\[
g(\lambda)
=
-\vec{b}^T (A-\lambda I)^\dagger \vec{b}
+\lambda.
\]

Thus the dual problem is

\[
d^\star =
\max_{\lambda\ge0}
\;
-\vec{b}^T (A-\lambda I)^\dagger \vec{b} + \lambda
\]

subject to

\[
A-\lambda I\succeq0, \qquad \vec{b}\in R(A-\lambda I).
\]

%----------------------------------------------------------------------------------------

\newpage
\section*{Problem 3. Minimizing a Sum of Logarithms}

Consider the following problem, which arises in estimation of transition probabilities of a discrete-time Markov chain:

\[
p^\star = \max_{\vec{x}\in\mathbb{R}^n} \sum_{i=1}^n \alpha_i \log(x_i)
\]

\[
\text{s.t. } \vec{x} \ge 0, \quad \mathbf{1}^T \vec{x} = c,
\]

where $c > 0$ and $\alpha_i > 0$, $i = 1, \dots, n$. (Recall that if $\vec{x}$ is a vector then by “$\vec{x} \ge 0$” we mean “$x_i \ge 0$ for each $i$.”) We will determine in closed-form a minimizer, and show that the optimal objective value of this problem is

\[
p^\star = \alpha \log(c/\alpha) + \sum_{i=1}^n \alpha_i \log(\alpha_i),
\]

where
\[
\alpha := \sum_{i=1}^n \alpha_i.
\]

We will show this in a series of steps.

\begin{enumerate}

\item[(a)] First, express the problem as a minimization problem which has optimal value $p^\star_{\min}$.

\item[(b)] In optimization, we often “relax” problems of the form 
\[
p^\star_{\min} = \min_{\vec{x}\in X} f_0(\vec{x}),
\]
i.e., replacing the constraint set $X$ with a larger constraint set $X_r$, and instead solving
\[
p^\star_r = \min_{\vec{x}\in X_r} f_0(\vec{x}),
\]
then showing a connection between $p^\star_{\min}$ and $p^\star_r$.

In this problem, a particular relaxation we will use is to replace the equality constraint $\mathbf{1}^T \vec{x} = c$ with an inequality constraint $\mathbf{1}^T \vec{x} \le c$.

Show that the relaxed problem has the same optimal value as the original problem, i.e., $p^\star_r = p^\star_{\min}$, and the two problems have the same solutions.

Hint: First argue that $p^\star_r \le p^\star_{\min}$. Then, suppose for the sake of contradiction that $p^\star_r < p^\star_{\min}$. Let $\vec{x}_r$ be a solution to the relaxed minimization problem which has objective value $p^\star_r$. Consider the vector $\vec{x}$ given by

\[
\vec{x} :=
\begin{bmatrix}
c - \mathbf{1}^T \vec{x}_r + x_{r1} \\
x_{r2} \\
\vdots \\
x_{rn}
\end{bmatrix}.
\]

Show that $\vec{x}$ is feasible for the original problem and has objective value $< p^\star_r$. Argue that this implies $p^\star_{\min} < p^\star_r$ and derive a contradiction. Finally, argue that any solution to the relaxed problem is a solution to the original problem, and vice-versa — you might need to use a construction similar to $\vec{x}$.

\item[(c)] After relaxing the equality constraint to an inequality constraint, form the Lagrangian $L(\vec{x}, \vec{\lambda}, \mu)$ for the relaxed minimization problem, where $\lambda_i$ is the dual variable corresponding to the inequality $x_i \ge 0$, and $\mu$ is the dual variable corresponding to the inequality constraint $\mathbf{1}^T \vec{x} \le c$.

\item[(d)] Now derive the dual function $g(\vec{\lambda}, \mu)$ for the relaxed minimization problem, and solve the dual problem

\[
d^\star_r = \max_{\vec{\lambda} \ge 0,\, \mu \ge 0} g(\vec{\lambda}, \mu).
\]

What are the optimal dual variables $\vec{\lambda}^\star, \mu^\star$?

\item[(e)] Show that strong duality holds for the relaxed problem, so $p^\star_r = d^\star_r$.

\item[(f)] From the $\vec{\lambda}^\star, \mu^\star$ obtained in the previous part, how do we obtain the optimal primal variable $\vec{x}^\star$? What is the optimal objective function value $p^\star_r$? Finally, what is $p^\star$?

\end{enumerate}

\subsection*{Answer}

\textbf{(a)}

\[
p^\star_{\min}
=
\min_{\vec{x}} -\sum_{i=1}^n \alpha_i \log(x_i)
\]

subject to

\[
\vec{x} \ge 0, \qquad \mathbf{1}^T\vec{x} = c.
\]

Thus

\[
p^\star_{\min} = -p^\star.
\]

\textbf{(b)}

Let the relaxed optimal value be $p_r^\star$.

Since the relaxed problem has a larger feasible set,

\[
p_r^\star \le p_{\min}^\star.
\]

Assume for contradiction that

\[
p_r^\star < p_{\min}^\star.
\]

Let $\vec{x}_r$ be the optimal solution of the relaxed problem.

Define

\[
\vec{x} =
\begin{bmatrix}
c - \mathbf{1}^T\vec{x}_r + x_{r1} \\
x_{r2} \\
\vdots \\
x_{rn}
\end{bmatrix}.
\]

Then

\[
\mathbf{1}^T\vec{x} = c,
\]

so $\vec{x}$ is feasible for the original problem.

Since $\log(x)$ is increasing, increasing one component increases the objective of the maximization problem (or decreases the minimization objective). Therefore

\[
-\sum_i \alpha_i\log(x_i) < -\sum_i \alpha_i\log(x_{ri}) = p_r^\star,
\]

which implies

\[
p_{\min}^\star < p_r^\star,
\]

a contradiction.

Thus

\[
p_r^\star = p_{\min}^\star,
\]

and both problems share the same solutions.

\bigskip

\textbf{(c)}

The relaxed minimization problem is

\[
\min_{\vec{x}} -\sum_{i=1}^n \alpha_i\log(x_i)
\]

subject to

\[
x_i \ge 0, \qquad \mathbf{1}^T\vec{x} \le c.
\]

The Lagrangian is

\[
L(\vec{x},\vec{\lambda},\mu)
=
-\sum_{i=1}^n \alpha_i\log(x_i)
-
\sum_{i=1}^n \lambda_i x_i
+
\mu(\mathbf{1}^T\vec{x}-c),
\]

where

\[
\lambda_i \ge 0, \qquad \mu \ge 0.
\]

\bigskip

\textbf{(d)}

The dual function is

\[
g(\vec{\lambda},\mu)
=
\inf_{\vec{x}} L(\vec{x},\vec{\lambda},\mu).
\]

Rearranging,

\[
L
=
-\sum_i \alpha_i\log(x_i)
+
\sum_i (\mu-\lambda_i)x_i
-
\mu c.
\]

Minimize with respect to $x_i$:

\[
\frac{\partial L}{\partial x_i}
=
-\frac{\alpha_i}{x_i} + \mu - \lambda_i = 0.
\]

Thus

\[
x_i = \frac{\alpha_i}{\mu - \lambda_i}.
\]

Substitute into $L$:

\[
g(\lambda,\mu)
=
-\sum_i \alpha_i\log\!\left(\frac{\alpha_i}{\mu-\lambda_i}\right)
+
\sum_i \alpha_i
-
\mu c.
\]

Maximizing this with $\lambda_i \ge 0$ yields the optimum at

\[
\lambda_i^\star = 0.
\]

Then

\[
x_i = \frac{\alpha_i}{\mu}.
\]

Using the constraint

\[
\mathbf{1}^T x = c,
\]

we obtain

\[
\frac{1}{\mu}\sum_i \alpha_i = c.
\]

Let

\[
\alpha = \sum_i \alpha_i.
\]

Thus

\[
\mu^\star = \frac{\alpha}{c}.
\]

\bigskip

\textbf{(e)}

The objective $-\sum \alpha_i\log(x_i)$ is convex and the constraints are linear.

Since there exists a strictly feasible point (e.g., $x_i = c/n >0$), Slater's condition holds.

Therefore strong duality holds:

\[
p_r^\star = d_r^\star.
\]

\bigskip

\textbf{(f)}

Using the optimal dual variables,

\[
x_i^\star = \frac{\alpha_i}{\mu^\star}
= \frac{c\alpha_i}{\alpha}.
\]

\[
p^\star
=
\sum_i \alpha_i
\log\!\left(\frac{c\alpha_i}{\alpha}\right).
\]

Expanding,

\[
p^\star
=
\sum_i \alpha_i
\left(
\log(c/\alpha) + \log(\alpha_i)
\right).
\]

Thus

\[
p^\star
=
\alpha\log(c/\alpha)
+
\sum_{i=1}^n \alpha_i\log(\alpha_i),
\]

where

\[
\alpha = \sum_i \alpha_i.
\]

This matches the stated optimal objective value.
%----------------------------------------------------------------------------------------

\newpage
\section*{Problem 4. Lagrangian Dual of a QP}


Consider the convex quadratic program

\[
\min_{\vec{x}} \frac{1}{2}\vec{x}^TQ\vec{x}
\]

subject to

\[
A\vec{x} \le \vec{b}
\]

where $Q\succ0$.

\begin{enumerate}
\item[(a)] Write the Lagrangian $L(\vec{x},\vec{\lambda})$.

\item[(b)] Write the Lagrangian dual function $g(\vec{\lambda})$.

\item[(c)] Show the dual problem is convex by writing it in standard QP form. Is the Lagrangian dual problem convex in general?
\end{enumerate}

\subsection*{Answer}

\textbf{(a)}

\[
L(\vec{x},\vec{\lambda})
=
\frac{1}{2}\vec{x}^T Q \vec{x}
+
\vec{\lambda}^T(A\vec{x}-\vec{b}).
\]

\bigskip

\textbf{(b)}

The dual function is

\[
g(\vec{\lambda}) = \inf_{\vec{x}} L(\vec{x},\vec{\lambda}).
\]

First rewrite the Lagrangian:

\[
L(\vec{x},\vec{\lambda})
=
\frac{1}{2}\vec{x}^T Q \vec{x}
+
\vec{\lambda}^T A \vec{x}
-
\vec{\lambda}^T \vec{b}.
\]

Minimize with respect to $\vec{x}$.  
Take the gradient:

\[
\nabla_x L = Q\vec{x} + A^T\vec{\lambda}.
\]

Set to zero:

\[
Q\vec{x} + A^T\vec{\lambda} = 0.
\]

Since $Q \succ 0$, it is invertible:

\[
\vec{x}^\star = -Q^{-1}A^T\vec{\lambda}.
\]

Substitute into the Lagrangian:

\[
g(\vec{\lambda})
=
\frac{1}{2}(-Q^{-1}A^T\vec{\lambda})^T
Q
(-Q^{-1}A^T\vec{\lambda})
+
\vec{\lambda}^T A (-Q^{-1}A^T\vec{\lambda})
-
\vec{\lambda}^T\vec{b}.
\]

Simplifying,

\[
g(\vec{\lambda})
=
-\frac{1}{2}\vec{\lambda}^T A Q^{-1} A^T \vec{\lambda}
-
\vec{\lambda}^T \vec{b}.
\]

Thus the dual problem is

\[
\max_{\vec{\lambda} \ge 0}
\;
-\frac{1}{2}\vec{\lambda}^T A Q^{-1} A^T \vec{\lambda}
-
\vec{\lambda}^T \vec{b}.
\]

\bigskip

\textbf{(c)}

The dual problem can be written equivalently as a minimization problem:

\[
\min_{\vec{\lambda} \ge 0}
\frac{1}{2}\vec{\lambda}^T A Q^{-1} A^T \vec{\lambda}
+
\vec{\lambda}^T \vec{b}.
\]

This has the standard quadratic program form

\[
\min_{\vec{\lambda}}
\frac{1}{2}\vec{\lambda}^T P \vec{\lambda} + \vec{b}^T\vec{\lambda}
\]

subject to

\[
\vec{\lambda} \ge 0,
\]

where

\[
P = A Q^{-1} A^T.
\]

Since $Q \succ 0$, we have $Q^{-1} \succ 0$, and therefore

\[
A Q^{-1} A^T \succeq 0.
\]

Thus the Hessian matrix $H$ is positive semidefinite, which means the dual problem is a convex quadratic program.

In general, the Lagrangian dual problem is always convex, even if the primal problem is not convex.

%----------------------------------------------------------------------------------------

\newpage
\section*{Problem 5. Dual of the Dual of a Linear Program}


Consider the linear program

\[
\min_{\vec{x}} \vec{c}^T\vec{x}
\]

subject to

\[
A\vec{x}= \vec{b},\quad \vec{x}\ge0.
\]

\begin{enumerate}
\item[(a)] Formulate the Lagrangian and write the dual problem (without $\vec{x}$).

\item[(b)] Express the dual as an equivalent minimization problem. Find the dual of this minimization problem (dual of the dual) and compare it to the original LP.
\end{enumerate}

\subsection*{Answer}

\textbf{(a)}

The primal linear program is

\[
\min_{\vec{x}} \vec{c}^T\vec{x}
\]

subject to

\[
A\vec{x} = \vec{b}, \qquad \vec{x} \ge 0.
\]

Introduce dual variables:
\[
\vec{y} \in \mathbb{R}^m \quad \text{for } A\vec{x}=b,
\]
\[
\vec{\lambda} \ge 0 \quad \text{for } x \ge 0.
\]

The Lagrangian is

\[
L(\vec{x},\vec{y},\vec{\lambda})
=
\vec{c}^T\vec{x}
+
\vec{y}^T(A\vec{x}-\vec{b})
-
\vec{\lambda}^T\vec{x}.
\]

Rearranging,

\[
L
=
(c + A^T y - \lambda)^T x - y^T b.
\]

The dual function is

\[
g(y,\lambda) = \inf_x L(x,y,\lambda).
\]

If

\[
c + A^T y - \lambda \neq 0
\]

then the problem is unbounded. Therefore we require

\[
c + A^T y - \lambda = 0.
\]

Thus

\[
\lambda = c + A^T y.
\]

Since $\lambda \ge 0$,

\[
c + A^T y \ge 0.
\]

The dual function becomes

\[
g(y) = -b^T y.
\]

Thus the dual problem is

\[
\max_y -b^T y
\]

subject to

\[
A^T y + c \ge 0.
\]

Equivalently,

\[
\max_y \; b^T(-y)
\quad
\text{s.t. } A^T(-y) \le c
\]

which is the standard LP dual form.

\textbf{(b)}

Rewrite the dual as an equivalent minimization problem by negating the objective:

\[
\min_y b^T y
\]

subject to

\[
A^T y + c \ge 0.
\]

Introduce dual variables $\vec{x} \ge 0$ corresponding to the constraint

\[
A^T y + c \ge 0.
\]

The Lagrangian is

\[
L(y,x)
=
b^T y + x^T(-A^T y - c).
\]

Rearranging,

\[
L
=
(b - A x)^T y - c^T x.
\]

The dual function is finite only if

\[
b - A x = 0.
\]

Thus

\[
Ax = b.
\]

Substituting,

\[
g(x) = -c^T x.
\]

Therefore the dual of the dual is

\[
\max_x -c^T x
\]

subject to

\[
Ax = b, \quad x \ge 0.
\]

Multiplying the objective by $-1$ gives

\[
\min_x c^T x
\]

subject to

\[
Ax = b, \quad x \ge 0.
\]

This is exactly the original primal linear program.

Thus, the dual of the dual of a linear program is the original linear program.

%----------------------------------------------------------------------------------------

\section*{Problem 6. Homework Process}


With whom did you work on this homework? List the names and SIDs of your group members.

NOTE: If you didn’t work with anyone, you can put “none”.

\subsection*{Answer}
none
%----------------------------------------------------------------------------------------

\end{document}
